\section{A directly certified dominant Gaussian $E_8$ component}
\label{sec:im-e8-period}

\paragraph{Provenance and precise assertion.}
This section proves that at least one twelve-branch component for
$\mu_4W(E_8)$ has period rank nine and raw Hermitian determinant
class $[-3]$. The existence proof uses a candidate point from Benjamin
Severn's unrefereed manuscript
\emph{Weil classes on abelian sixfolds of Weil type for
$\mathbf Q(i)$ with discriminant $-3$}, draft 0.7, 21 October 2026,
Hexagon 2610.00477v1, Appendix B, and its numerical archive
Zenodo 23549514. A supplied approximate point is not a theorem input.
The exact target matrices, 87 rational candidate coordinates, and
nine fixed gauge coordinates are retained as mathematical input;
the contraction, its side conditions, and the period minor are
recomputed. The source's approximate continuation, its claimed
outputs, and its main theorems are not premises.

The programs \texttt{e8\_ball/verify\_M.py} and
\texttt{e8\_ball/verify\_L.py} are explicitly acknowledged adaptations
of the source algorithms, whose complete implementations and error
bounds were read and checked. They replace the unavailable continuation
input by the saved candidate, use safe JSON data, and recompute every
certificate inequality. The target-group and path program
\texttt{check\_group\_and\_paths.py} was independently written.
The proof below supplies the mathematical content of the certification.
It proves dominance for the component containing this smooth cover;
it does not identify that component with the two-cancelling-pair
component D-E8T-01.

\subsection*{D-E8R-01. Exact target monodromy and the covering construction}
Let $H$ be the Cartan matrix of D-E8T-01. In the fundamental-weight
basis, whose Gram matrix is $H^{-1}$, use
\[
a=\begin{pmatrix}
0&0&1&0&0&0&0&0\\
1&1&1&1&1&1&0&0\\
1&1&1&2&2&2&2&1\\
-1&-1&-1&-1&-1&-1&-1&0\\
-1&-1&-2&-3&-3&-2&-1&-1\\
0&0&0&0&1&0&0&0\\
1&1&2&3&2&2&1&0\\
1&2&2&3&2&1&1&1
\end{pmatrix},
\]
\[
b=\begin{pmatrix}
-1&-2&-3&-5&-4&-3&-2&-1\\
2&3&3&5&4&3&2&1\\
2&2&4&5&4&3&2&1\\
-2&-2&-3&-4&-4&-3&-2&-1\\
0&0&0&0&1&0&0&0\\
0&0&0&0&0&1&0&0\\
0&0&0&0&0&0&1&0\\
0&0&0&0&0&0&0&1
\end{pmatrix},\qquad
c=I-e_1e_1^t+(-e_1+e_3)e_1^t.
\]
Thus $c$ is the simple reflection in the first root in this basis.
The target input \texttt{mtuple.json} also records $d=(abc)^{-1}$,
although identifying the order of a product of based loops will not
be needed for the analytic argument.

The three matrices belong to $W(E_8)$ and generate it. Their exact
membership certificate successively right-multiplies
$H^{-1}aH,H^{-1}bH,H^{-1}cH$ by simple reflections until the identity
is reached, retaining all word indices. Starting with the first
simple root and closing under these three matrices in the simple-root
basis gives all 240 roots. Since $c$ is a root reflection, conjugation
then supplies every root reflection, proving generation. The same
certificate checks $\chi_a=\Phi_{20}$ and that $b,c$ are involutions
with one-dimensional moving spaces. Thus no order-only identification
of the generated group is involved.

We will certify a Fuchsian system
\[
\frac{dY}{dz}=
 \left(\frac{B_0}{z}+\frac{B_1}{z-1}
                    +\frac{B_2}{z-(2+i)}\right)Y,
\qquad
B_0=\operatorname{diag}(1,3,7,9,11,13,17,19)/20,
\tag{R1}
\]
whose three based monodromies are simultaneously conjugate to
$a,b,c$. Here $B_1=u_1v_1^t$, $B_2=u_2v_2^t$ have rank one and
trace $1/2$. Pull (R1) back by $z=s^5$ and tensor its solution with
$s^{1/4}$. Its monodromy is the full group $G=\mu_4W(E_8)$.

To prove the latter statement without assuming a particular formula
for a twelve-entry based tuple, let $\Pi$ be the fundamental group
of the four-punctured $z$-sphere and let $\Pi'$ be the subgroup for
$s^5$. It is the kernel of the epimorphism to $\mathbf Z/5$ which
sends the loop around zero to one, those around $1,2+i$ to zero, and
that around infinity to minus one. The image of $\Pi'$ in $W(E_8)$
is normal with cyclic quotient of order dividing five. Since a group
generated by involutions has no nontrivial quotient of order five,
this image is all of $W(E_8)$. After the scalar twist the projective
image remains $W(E_8)/\{\pm I\}$. At each of the five preimages of
one, the twisted system still has a real root reflection as local
monodromy. Conjugating one such reflection by the projective image
supplies all root reflections, hence all of $W(E_8)$. The monodromy
at zero is conjugate to $ia^5$, which is outside $W(E_8)$. Therefore
the full image is $\mu_4W(E_8)$.

The finite monodromy kernel supplies, by Riemann existence, a smooth
connected algebraic $G$-cover of the $s$-sphere. Its twelve branch
points are zero, the fifth roots of one, the fifth roots of $2+i$,
and infinity. The exact local exponents and the Weil signature are
verified below after the existence certificate.

\subsection*{D-E8R-02. Rigorous integration and a finite contraction test}
The candidate parameters are the entries of
\[
x=(v_1,\ (u_{2,j})_{j\ne j_0},\ v_2,\ Q)\in\mathbf C^{87}.
\]
The eight entries of $u_1$, the index $j_0=1$ in zero-based notation,
and the coordinate $u_{2,j_0}$ are fixed dyadic rational complex
numbers, recorded in \texttt{center\_g.json}. The 87 candidate
coordinates are likewise specified exactly as pairs of integers
$(m,e)$ representing $m2^e$ at precision 256. The decimal source
strings are retained only to make extraction auditable. The program
asserts that their 256-bit midpoints equal the recorded dyadics.

We use the maximum absolute-value norm on vectors and its induced
maximum-row-sum matrix norm. A complex interval is a rectangular
product of two real Arb intervals. Each binary floating-point number
selected by NumPy is thereafter treated as its exact dyadic rational
value. Floating linear algebra chooses pivots or preconditioners;
it never certifies invertibility or decides whether a period minor
is nonzero. All such decisions use enclosing Arb calculations.

\paragraph{Transport error bound.}
For a regular expansion center $c_0$ of an equation
$Y'=A(z)Y$, $A(z)=\sum_k A_k/(z-p_k)$, suppose
\[
d\le\min_k|c_0-p_k|,\qquad S\ge\sum_k\|A_k\|_\infty.
\]
The Taylor coefficients $A_n$ at $c_0$ satisfy
$\|A_n\|_\infty\le S/d^{n+1}$. If
$T(z)=\sum_{n\ge0}T_n(z-c_0)^n$ is the transport with $T_0=I$,
its recursion is
\[
(n+1)T_{n+1}=\sum_{j=0}^n A_jT_{n-j}.
\]
Coefficientwise induction compares its norms with the scalar solution
of $f'=S(d-t)^{-1}f$, $f(0)=1$. Consequently
\[
\|T_n\|_\infty\le\frac{(S)_n}{n!d^n}.
\]
For a step $h$, put $q\ge|h|/d$, with $q<1/2$. After degree $N$,
the absolute tail in each matrix entry is at most
\[
\epsilon_N=
\frac{(S)_{N+1}q^{N+1}/(N+1)!}
 {1-q(1+S/(N+2))},
\tag{R2}
\]
provided the denominator is positive. Indeed the ratio of consecutive
terms after the first omitted term is at most
$q(1+S/(N+2))$. This argument is uniform when the residues range
over an input interval box. Adding $\epsilon_N$ to both coordinate
radii encloses the omitted complex entry.

The code computes either this Taylor recursion or the equivalent
polynomial recurrence obtained by clearing
$\prod_k(z-p_k)$. It evaluates the truncated transport by Horner's
rule and multiplies successive interval matrices. The two recurrences
solve the same differential equation and use exactly the same bound
(R2); the direct Taylor version is used for wider input boxes to
reduce interval inflation. Bounds exported to binary floats use the
appropriate Arb endpoint followed by one \texttt{nextafter} step
outward. Products and sums used in certificate inequalities are
returned to Arb before they are compared. No printed rounded decimal
is used as an inequality test.

The three polygonal loops start at $1/2-i$, travel to the near side
of the radius-$0.45$ circle about the relevant pole, traverse a
counterclockwise twelve-gon, and return. Their vertices are exact
binary rationals after selection. The independently written path
check verifies winding numbers $(1,0,0),(0,1,0),(0,0,1)$ by exact
rational cross products; after shrinking the small circles to radius $1/10$, their straight
access paths are disjoint except at the base point and avoid the other
small discs, giving the usual free generating loops. The exact minimum
distance from any other pole to such a straight access segment is
$1/5$; the polygon vertices and their access segments are homotopic
to these loops within the corresponding puncture complements. Adaptive steps
are dyadic and are checked to have $|h|/d<1/2$. All endpoint
additions use at most 56 numerator bits, hence are exact at precision
256. The actual three paths use 31, 31, and 53 steps. The final step
is the exact difference to the endpoint; an explicit assertion rules
out premature termination caused by a rounded endpoint.
There is also an independent exact rational comparison of the rounded
adaptive polyline with each intended polygon edge, implemented in
\texttt{e8\_path\_tube\_certificate.py}. Every adaptive vertex is
within $10^{-12}$ of that edge, while every intended edge is at
distance greater than $1/10$ from every pole. The closed
$10^{-12}$-neighbourhood of a line segment is convex and contains
no pole. Both its adaptive subpath and its intended segment therefore
lie in this neighbourhood and are homotopic relative to their common
endpoints. Concatenation gives the claimed based-loop homotopy.
Thus the computed matrices are monodromies of the specified closed
generating loops, not merely of paths with the same winding numbers.

\paragraph{The 87 equations and their side conditions.}
Let $M_0,M_1,M_2$ be these monodromies as holomorphic functions of
the residue parameters. For the eight distinct eigenvalues
$\lambda_j=\exp(2\pi i m_j/20)$ of $a$, construct exact algebraic
eigenvectors $e_j$ using a seven-row solve and one coordinate equal
to one. Ball inversion encloses those eigenvectors. Write
$b=I-x_by_b^t$ and $c=I-x_cy_c^t$ exactly, with
$y_b^tx_b=y_c^tx_c=2$, and choose rational bases of
$\ker y_b^t$ and $\ker y_c^t$.
The equations $F(x)=0$ are the following:
\begin{enumerate}
\item $v_1^tu_1-1/2=0$ and $v_2^tu_2-1/2=0$;
\item seven selected rows of $(M_0-\lambda_jI)Qe_j=0$ for each $j$;
\item seven selected rows of $(M_1+I)Qx_b=0$, and one fixed row
of $(M_1-I)Qw=0$ for each of seven basis vectors $w\in\ker y_b^t$;
\item the analogous fourteen equations for $M_2,c$;
\item $\det Q-1=0$.
\end{enumerate}
Their number is $2+56+14+14+1=87$. The dropped rows for the eight
$M_0$ equations are $0,1,\ldots,7$. The negative-eigenvector dropped
row and positive-eigenvector pivot row are $(4,3)$ for $M_1$ and
$(3,1)$ for $M_2$, in zero-based notation.

At a zero, the rank-one residues have trace $1/2$, so their local
monodromies are reflections. The fixed residue $B_0$ has distinct
nonresonant exponents and its monodromy has the stated eight simple
eigenvalues. The following additional inequalities are certified
throughout the small solution ball: the seven retained rows of
$M_0-\lambda_jI$ have rank seven; the seven retained rows of each
$M_k+I$, $k=1,2$, have rank seven; and the selected row of $M_k-I$
is nonzero. Each rank assertion is witnessed by a $7\times7$
minor excluding zero. Hence the seven-row kernels in the equations
are the full kernels. For $M_k-I$, its rank is one at a solution,
so the kernel of a nonzero row is its full kernel. Thus a zero
satisfies
\[
M_0Q=Qa,\qquad M_1Q=Qb,\qquad M_2Q=Qc,\qquad \det Q=1.
\tag{R3}
\]
This justifies the reduction to 87 equations; none of the omitted
rows is inferred merely from a small numerical residual.

\paragraph{Proof of existence.}
The map $F$ is holomorphic in all 87 parameters: fundamental-solution
transport along fixed nonsingular paths depends holomorphically on
the residues. Let $x_0$ be the recorded dyadic candidate, $n=87$,
$R=2^{-40}$, and $h=2^{-80}$. Compute an interval enclosure of
$F(x_0)$ and of every forward difference
\[
D_{ij}=\frac{F_i(x_0+he_j)-F_i(x_0)}h.
\]
For the 64 variables belonging only to $Q$, the same center monodromy
matrices can be reused exactly; this removes redundant ODE solves
without changing the forward differences. Let $Y_0$ be the exact
dyadic matrix obtained from a floating approximate inverse of the
midpoint of $D$. Enclose $F$ on the complex polydisc of radius $R$
by evaluating it on the containing real-coordinate interval box,
and write $K_i$ for an upper bound for $|F_i|$ there.

Cauchy's estimate, applied on the $j$th coordinate disc, gives
\[
|D_{ij}-\partial_jF_i(x_0)|
 \le\frac{K_i h}{R(R-h)}.
\tag{R4}
\]
For completeness, the coefficients of the one-variable Taylor
series are bounded by $K_i/R^k$; subtracting its first derivative
from the forward difference leaves
$\sum_{k\ge2}K_i h^{k-1}/R^k$, which is (R4).
For $\|x-x_0\|_\infty\le\rho<R$, twice applying Cauchy's estimate
and integrating along the segment from $x_0$ to $x$ gives
\[
|\partial_jF_i(x)-\partial_jF_i(x_0)|
 \le\frac{2n\rho K_i}{(R-\rho)^2}.
\tag{R5}
\]
The factor two includes the equal-coordinate second derivative;
for distinct coordinates the smaller factor one would suffice.
Summing a row and using submultiplicativity therefore gives
\[
\begin{split}
\eta&=\|Y_0F(x_0)\|_\infty,\\
E&=\max_i\frac{nK_i h}{R(R-h)},\qquad
C=\max_i\frac{2n^2\rho K_i}{(R-\rho)^2},\\
\kappa&=\|I-Y_0D\|_\infty+\|Y_0\|_\infty(E+C).
\end{split}
\tag{R6}
\]
In the actual computation $\rho$ is the outward-rounded value
$10\eta$. The fresh run gives, with these figures rounded only for
display,
\[
\eta\simeq1.052812\,10^{-35},\quad
\rho\simeq1.052812\,10^{-34},\quad
\kappa\simeq0.005589859013,
\]
\[
\max_i K_i<1.172\,10^{-6},\qquad
\|I-Y_0D\|_\infty<2.216\,10^{-9},\qquad
\|Y_0\|_\infty<54.87.
\]
Arb verifies $\kappa<1$ and $\eta+\kappa\rho<\rho$.
The map $x\mapsto x-Y_0F(x)$ is consequently a strict contraction
of the closed maximum-norm ball of radius $\rho$ about $x_0$ into
itself. It has a unique fixed point. Moreover
$\|I-Y_0D\|_\infty<1$ makes $Y_0D$ invertible and therefore makes
$Y_0$ invertible, so a fixed point has $F(x)=0$. The side-minor tests
are performed on the entire enclosing radius-$\rho$ input box.
This proves existence of (R1) with the exact finite target monodromy
(R3), not merely proximity to finite monodromy.

\subsection*{D-E8R-03. Exact exponents and the polarized Weil factor}
Put $B_\infty=-B_0-B_1-B_2$. A numerical eigenbasis of its midpoint
is treated as an exact dyadic invertible matrix. Interval inversion
and multiplication then give Gershgorin discs, each contained within
$10^{-10}$ of the following rational labels, counted with
multiplicity:
\[
-1/12,\ -1,\ -11/12,\ -2/3,\ -7/12,\ -1/3,\ -5/12,\ -1.
\tag{R7}
\]
This localization is recomputed in Arb in both verification stages;
the displayed floating disc radii in the diagnostic log are not used
as a proof. The disjoint clusters of discs count one eigenvalue at
each single label and two at $-1$.

The labels are exact. By the residue--monodromy theorem, the
exponential of $2\pi i$ times a residue eigenvalue is a local
monodromy eigenvalue; see \cite[Corollary 16.20, author extract
p.~274]{IM:IlyashenkoYakovenko}. Every monodromy here is conjugate
to an element of the rational eight-dimensional representation of
$W(E_8)$. An eigenvalue of primitive order $e$ therefore has
$\varphi(e)\le8$. The complete possible list is
\[
1,2,3,4,5,6,7,8,9,10,12,14,15,16,18,20,24,30.
\]
One can verify completeness without a class table: each prime
$p\mid e$ has $p-1\le8$, hence $p\in\{2,3,5,7\}$; their possible
powers divide $2^4,3^2,5,7$, and checking their products gives exactly
the displayed list. Different candidates have rational angles
separated by far more than $10^{-10}$ from the labels (R7).
Thus the residue eigenvalues are precisely (R7), even without
fixing an ordering convention for the based-loop product at infinity.

For later rational differential calculations use the integer scalar
gauge $s^{-3}$ in addition to the scalar twist. The local exponents
of the pulled-back system are
\[
5m_j/20+1/4-3\quad\text{at }s=0,
\qquad (0^7,1/2)\quad\text{at each of the ten reflection points},
\]
\[
5e_j-1/4+3\quad\text{at }s=\infty,
\qquad e_j\text{ as in (R7)}.
\tag{R8}
\]
The gauge changes no monodromy. At zero these exponents modulo one
are four zeros and four halves. At infinity they are
\[
1/12,\ 1/6,\ 1/3,\ 5/12,\ 2/3,\ 3/4,\ 3/4,\ 5/6.
\]
Their sum is four. Thus the moving-rank and age sums are 22 and 11,
respectively. Chevalley--Weil, in the form
\cite[Theorem 2.10, author manuscript p.~8]{IM:FredianiGhigiPenegini},
gives the two holomorphic multiplicities $11-8=3$.
The natural character has field $K=\mathbf Q(i)$ and is realizable
over $K$, so a rational Rosati-fixed primitive projector gives an
actual polarized abelian sixfold $A_0$, with signature $(3,3)$.
This is the same group-algebra construction used in D-E8T-02.

Its determinant class is $[-3]$. Every real reflection has local
factor one in the determinant formula D-DET-01, using a root of
norm two. For the element $ia^5$, the first independent columns of
$I-ia^5$ give a moving-space Gram determinant 16 and determinant
$\det(I-ia^5|_{\mathrm{mov}})=16$, so its factor is one. At infinity
all eight directions move, the full Gram determinant is one, and
the displayed spectrum gives $\det(I-g_\infty)=6$; its local factor
is $1/6$. The signed six-dimensional determinant class is therefore
$[1/6]=[3]$, and the raw determinant class is $[-3]$. These are
computations for the actual invariant positive form, not an arbitrary
auxiliary Hermitian form.

\subsection*{D-E8R-04. Rational descriptions of the Hodge eigenspaces}
Change coordinates to $t=1/(s-s_*)$, where $s_*$ is the exact
binary rational complex value represented by $0.37+0.21i$. Then
$t=\infty$ is a regular point. The system has twelve finite simple
poles $p$, with residues $A_p$, whose exponents are (R8), and total
residue zero. Let $Y$ be a fundamental solution. Let $D_+$ and
$D_-$ denote the two three-dimensional natural multiplicity spaces
in $H^0(C_0,\omega_{C_0})$.

There are canonical identifications, up to interchanging the two
embeddings of $K$,
\[
D_+=\{r\in\mathbf C(t)^8:Y^{-1}r\,dt
                        \text{ is holomorphic on }C_0\},
\]
\[
D_-=\{q\in\mathbf C(t)^8:Y^tq\,dt
                        \text{ is holomorphic on }C_0\}.
\tag{R9}
\]
Indeed $Y$ becomes single-valued after pullback to the finite
monodromy cover. Its entries and those of its inverse have moderate
growth at the branch points, and hence extend meromorphically on
that cover. If $r$ is rational, $Y^{-1}r\,dt$ has the appropriate
contragredient transformation law. Conversely, for a vector of
holomorphic forms with that transformation law, multiplying by $Y$
and dividing by $dt$ gives a $G$-invariant meromorphic vector on the
cover. It descends to a meromorphic vector on the sphere and is
therefore rational. The argument for $Y^tq$ is identical with the
dual representation. These are spaces of algebraic holomorphic
forms on the actual compact algebraic curve, not spaces defined only
by a numerical differential equation.

Choose a residue eigenbasis $P_p$ and write a local fundamental
solution in the form
\[
Y=P_pH_p(t-p)(t-p)^{\Lambda_p},\qquad H_p(0)=I,
\qquad \Lambda_p=\operatorname{diag}(\lambda_{p,j}).
\tag{R10}
\]
Such a holomorphic frame exists also at resonant poles in this
calculation. One justification from the published local theory is
as follows. The holomorphic Poincare--Dulac--Levelt form has
monodromy $\exp(2\pi i\Lambda)\exp(2\pi iN)$, with commuting
semisimple and unipotent factors, where $N$ is nilpotent
\cite[Theorems 16.15--16.16 and Corollary 16.19,
author extract pp.~270--274]{IM:IlyashenkoYakovenko}.
Finite monodromy forces $\exp(2\pi iN)=I$, hence $N=0$.
The resonant coefficients and the residue's nilpotent part occupy
disjoint matrix-entry positions according to the eigenvalue
difference, so each vanishes. The normal form is the Euler system,
which gives (R10). Alternatively, before the fifth pullback all
four residue spectra lie in intervals of length less than one;
nonresonant Euler frames pull back to (R10).

The holomorphy conditions in (R9) are exactly
\[
\operatorname{ord}_p[(P_pH_p)^{-1}r]_j\ge\lfloor\lambda_{p,j}\rfloor,
\qquad
\operatorname{ord}_p[(P_pH_p)^tq]_j\ge\lfloor-\lambda_{p,j}\rfloor,
\tag{R11}
\]
together with regularity away from these poles and vanishing to
order at least two at infinity. To verify the floor precisely, let
$t-p=y^e$ on the cover. The first form's $j$th component has order
\[
e\operatorname{ord}_p[(P_pH_p)^{-1}r]_j
                    -e\lambda_{p,j}+e-1.
\]
As $e\lambda_{p,j}$ is an integer, this is nonnegative exactly when
the first inequality (R11) holds. The second follows by replacing
$\lambda$ by $-\lambda$. At infinity $dt$ has a pole of order two
and the system is regular, giving the stated vanishing condition.

Here is why the resonant recursion used in the certificate computes
a valid frame. Write $H_p=I+\sum_{m\ge1}H_mt^m$. Its $(i,j)$
coefficient at order $m$ has denominator
$m+\lambda_j-\lambda_i$. At a zero denominator choose the free
coefficient $(H_m)_{ij}$ to be zero. Starting from any frame (R10),
right multiplication of its fundamental matrix by a constant
unipotent matrix changes $H$ to
$H(t)t^\Lambda C t^{-\Lambda}$. An entry $C_{ij}$ for
$\lambda_i-\lambda_j=m>0$ first changes the resonant coefficient
of degree $m$. Choosing these entries in increasing order of $m$
sets all such coefficients to zero without altering lower ones.
Hence this normalization exists. The remaining recursion coefficients
are uniquely forced and converge by the holomorphic classification
theorem. The exact resonant right-hand sides are zero for this
geometric reason; the numerical test that their intervals contain
zero is only a consistency check, not a proof of vanishing.

\subsection*{D-E8R-05. Finite linear algebra for the actual codifferential}
For the first space in (R9), put
$M_p=\max(0,-\min_j\lfloor\lambda_{p,j}\rfloor)$, and use the
analogous definition with $-\lambda$ for the second. A rational
vector in either space has the unique partial-fraction expression
\[
r(t)=\sum_p\sum_{m=1}^{M_p}c_{p,m}(t-p)^{-m}.
\tag{R12}
\]
This follows because $P_pH_p$ and its inverse are holomorphic,
so these are upper bounds on all pole orders; vanishing at infinity
eliminates a polynomial part. Its order-two vanishing at infinity
adds the eight linear equations $\sum_p c_{p,1}=0$.

Expanding (R11) imposes finitely many linear equations on the
coefficients (R12). For the present exponents they give matrices
\[
\Phi_+:\mathbf C^{48}\longrightarrow\mathbf C^{45},\qquad
\Phi_-:\mathbf C^{120}\longrightarrow\mathbf C^{117}.
\]
The code computes the local frames through order twelve. Every
coefficient required by (R11) is within this range, as is checked
from the exact floors and pole orders; no infinite series tail is
used for this step. The eigenbases themselves are enclosed by
solving exact-eigenvalue kernel equations using square pivot blocks.
Their true kernel dimensions are known from the already proved
semisimplicity and exact exponent multiplicities.

Select 45 dependent columns in $\Phi_+$ and 117 in $\Phi_-$.
The corresponding square blocks $S_+,S_-$ have certified inverse
enclosures. More explicitly, using exact midpoint inverse
preconditioners $R_+,R_-$, the fresh run bounds
\[
\|I-R_+S_+\|_\infty<1.650\,10^{-24},\qquad
\|I-R_-S_-\|_\infty<3.385\,10^{-26}.
\]
Therefore both blocks are invertible. Set each of the remaining
three free coordinates in turn to one, the others to zero, and solve
the selected rows. These vectors form the kernel basis of those
rows. They are the complete kernel bases of $\Phi_\pm$: by
Chevalley--Weil the true kernels have dimension three, they lie
inside the three-dimensional selected-row kernels, and hence equal
them. Checking that all unused-row intervals contain zero is a
consistency test; the dimension argument, not that interval test,
proves that all rows vanish exactly.

We next prove that the resulting matrix is the period
codifferential. For a deformation of a curve, represent its
Kodaira--Spencer class by vector fields $v_{ab}$ on overlaps of
coordinate charts. Differentiating the gluing of a holomorphic
one-form $\omega$ gives $\mathcal L_{v_{ab}}\omega
=d(\iota_{v_{ab}}\omega)$. Thus the infinitesimal variation class
is represented in $H^1(\mathcal O_C)$ by
$\iota_{v_{ab}}\omega$. Pairing with a second holomorphic one-form
$\eta$ by Serre duality pairs the Kodaira--Spencer class with the
quadratic differential $\omega\eta$. This is the multiplication
description of the Torelli codifferential, derived here on charts.

For fixed-$G$ deformations, take invariants. Locally a cover has
$t=y^e$. An invariant vector field upstairs is
$y f(y^e)\partial_y$, whose image is
$e t f(t)\partial_t$. Hence
$(\pi_*T_C)^G=T_{\mathbf P^1}(-B)$, and the Hurwitz tangent space
is $H^1(T_{\mathbf P^1}(-B))$. Dually, invariant quadratic
differentials upstairs identify with
$H^0(\mathbf P^1,\omega_{\mathbf P^1}^{\otimes2}(B))$.
Indeed a possible pole downstairs has pullback order
$e\operatorname{ord}(q)+2e-2$, which for $e\ge2$ permits at most
a simple pole. On the natural nonreal representation pair, Schur's
lemma identifies the invariant contraction with the tensor product
of the two multiplicity spaces. Restricting the multiplication map
to this block is precisely the codifferential of the projected
abelian factor: the projector is an algebraic flat endomorphism of
the Jacobian variation. In the coordinates (R9) the contraction is
\[
\mu:D_+\otimes D_-\longrightarrow
 H^0(\mathbf P^1,\omega_{\mathbf P^1}^{\otimes2}(B)),\qquad
\mu(r\otimes q)=q^tr\,dt^2.
\tag{R13}
\]
Indeed $(Y^tq)^t(Y^{-1}r)=q^tr$. Conventions may multiply this
identification by a fixed nonzero scalar; they cannot change its rank.

All twelve branch points are finite in the $t$ coordinate. A
quadratic differential in the target is uniquely determined by its
simple-pole coefficients $\rho_p$, and these satisfy exactly
\[
\sum_p\rho_p=\sum_pp\rho_p=\sum_pp^2\rho_p=0.
\tag{R14}
\]
To see this, its rational coefficient is
$\sum_p\rho_p/(t-p)$, and regularity at infinity requires decay
$O(t^{-4})$. Expanding at infinity gives precisely (R14).
Thus the target has dimension $12-3=9$. For the three certified
bases on each side of (R13), form the nine rational contractions
$q_b^tr_a$, expand by the finite partial fractions (R12), and take
their twelve residues. The $9\times12$ residue matrix has the same
rank as the true codifferential.

\subsection*{D-E8R-06. A nonzero certified minor and dominance}
Using the freshly established solution enclosure, the minor on
zero-based residue columns $3,4,5,6,7,8,9,10,11$ is contained in
\[
[0.10557441259506563\ \mathord\pm\ 6.51\,10^{-18}]
+i[-0.04409991368479836\ \mathord\pm\ 3.88\,10^{-18}].
\tag{R15}
\]
The enclosure excludes zero, and its absolute-value lower bound is
greater than $0.11441485472524$. Matrix determinants, basis solves,
local frame coefficients, branch points, and residues here are all
evaluated in enclosing 256-bit arithmetic. The approximation of the
residues of (R1) is included from the outset as an input interval of
radius $\rho$ about the proven exact center. Consequently (R15)
contains the determinant at the exact finite-monodromy solution,
not just at its numerical midpoint. The three residue relations
(R14) are also checked as interval consistency conditions.

It follows that the differential of the actual polarized natural
factor period map at $C_0$ has rank nine. Its source Hurwitz
component has dimension nine, and the target signature-$(3,3)$
unitary domain has dimension nine. Passing to a finite marking and
level cover gives an actual family of abelian factors. The holomorphic
inverse-function theorem makes the local period image open in that
domain. The algebraic map to the fixed polarization and endomorphism
moduli component is therefore dominant. By D-E8R-03 that component
has raw determinant class $[-3]$.

This conclusion proves existence of a dominant nonsplit Gaussian
$E_8$ natural source. It gives neither a classification of its
Nielsen components nor a proof that every component, or the
particular two-pair component D-E8T-01, has rank nine. A supply of
algebraic Weil classes and a uniform bound for their detecting
cycles remain separate questions.

\paragraph{Reproduction and source limits.}
Run, in \texttt{agents/imprimitive/e8\_ball}, first
\texttt{python check\_group\_and\_paths.py}, then
\texttt{python verify\_M.py}, then \texttt{python verify\_L.py}.
The first step is exact rational algebra. The second is the
existence and side-condition proof; the third reads its newly
written enclosure and proves (R15). The successful run used
Python 3.12, python-flint 0.9.0, and 256-bit Arb arithmetic
\cite{IM:JohanssonArb}; NumPy and SciPy serve only for candidate
pivots and approximate preconditioners. The original continuation
solution is unnecessary and is not silently assumed to exist.
The compact input/provenance manifest records source and extracted
hashes. The full third-party archive is not part of the deliverable.
